\section{附录}

\begin{frame}{可能被问到的问题：识别与测度}
  \begin{description}
    \item[为什么不用年报文本词频来衡量数字化？] 词频指标有测度误差，而且和企业经营表现同期相关——说不清是数字化在影响投资，还是业绩波动同时影响了词频和投资。贯标的好处是有一个清晰、外生、可验证的处理时点。
    \item[为什么不用 KZ 或 SA 指数？] KZ 和 SA 度量的是融资约束的松紧程度，而我们关心的是：在控制了生产基本面之后，企业的资本存量有没有"多余"的部分。这是两个不同的问题。
    \item[会不会只是整体缩表？] R\&D 支出、研发人员占比、创新补助和研发资本化比率都在上升——这和整体收缩的画面对不上。
  \end{description}
\end{frame}

\begin{frame}{可能被问到的问题：模型与因果}
  \begin{description}
    \item[门槛的存在是否依赖于参数设定？] 门槛的存在来自 ABL 和 CFL 双契约结构本身，不是参数校准的结果。而且资本年龄异质性提供了独立于参数选择的经验支持。
    \item[会不会是企业自选择？] 事前企业特征对政策进入年份的解释力有限，PSM-DID、合成控制 DID 和堆叠 DID 的结论都和基准一致。
  \end{description}
\end{frame}

\begin{frame}{备用表：借款能力与融资约束}
  {\scriptsize
\begin{center}
\begin{tabular}{lcccc}
  \toprule
  Sun--Abraham 估计 & 总贷款 & 抵押贷款 & 短期贷款 & 长期贷款 \\
  \midrule
  两化融合 & $0.269^{***}$ & $0.210^{**}$ & $0.221^{***}$ & 0.148 \\
           & (0.068) & (0.100) & (0.069) & (0.094) \\
  \midrule
  融资约束缓解 & CapEX 全样本 & CapEX 民营 & 现金持有变动 &  \\
  \midrule
  经营现金流 & $0.023^{***}$ & $0.019^{***}$ & $0.356^{***}$ &  \\
             & (0.004) & (0.005) & (0.010) &  \\
  两化融合 $\times$ 现金流状况 & $-0.020^{*}$ & $-0.028^{**}$ & $-0.044^{*}$ &  \\
                                & (0.011) & (0.013) & (0.025) &  \\
  \bottomrule
\end{tabular}
\end{center}
}

\end{frame}
